\documentclass[../../main.tex]{subfiles}% 注意这里的文件路径不能用 ./main.tex ,否则用latexmk编译子文件会报错
\graphicspath{{\subfix{./image/}}{\subfix{../../image/}}} % 指定图片引用路径,后续可以直接使用图片文件名
% 如果图片存放在上级目录,则使用 ../../image/ 作为备用路径

% 例如:
% \begin{figure}[H]
% \centering
% \includegraphics[scale=0.3]{图.png}
% \caption{}
% \label{figure:图}
% \end{figure}
% 注意:上述\label{}一定要放在\caption{}之后,否则引用图片序号会只会显示??.

\begin{document}

\section{三维欧氏空间中的标架}

\begin{definition}[标架]
在 $E^3$ 中取定不共面的 $4$ 个点，把其中一点记作 $O$，把另外 $3$ 点分别记为 $A, B, C$，于是得到由一点 $O$ 和 $3$ 个不共面的向量 $\overrightarrow{OA}, \overrightarrow{OB}, \overrightarrow{OC}$ 构成的图形 $\{O; \overrightarrow{OA}, \overrightarrow{OB}, \overrightarrow{OC}\}$. 这样的一个图形称为 $E^3$ 中的一个\textbf{标架}，其中点 $O$ 称为该标架的\textbf{原点}.
\end{definition}

\begin{definition}[正交标架]
设 $\{O; \bm i, \bm j, \bm k\}$ 是 $E^3$ 的一个标架，并且 $\bm i, \bm j, \bm k$ 是彼此垂直的、构成右手系的 $3$ 个单位向量，于是
\begin{align*}
\bm i \cdot \bm i = \bm j \cdot \bm j = \bm k \cdot \bm k = 1,\quad \bm i \cdot \bm j = \bm i \cdot \bm k = \bm j \cdot \bm k = 0,
\end{align*}
并且
\begin{align*}
\bm i \times \bm j = \bm k,\quad \bm j \times \bm k = \bm i,\quad \bm k \times \bm i = \bm j.
\end{align*}
这样的标架称为\textbf{右手单位正交标架}，有时也简称为\textbf{正交标架}.
\end{definition}

\begin{definition}[等距变换和刚体运动]
设 $\sigma: E^3 \to E^3$ 是一个变换. 若对任意两点 $p,q\in E^3$ 都有
\begin{align*}
|\sigma(p)\sigma(q)| = |pq|,
\end{align*}
则称 $\sigma$ 为 $E^3$ 的一个\textbf{等距变换}. 若等距变换 $\sigma$ 还将右手正交标架映为右手正交标架 (即保持定向), 则称 $\sigma$ 为 $E^3$ 的一个\textbf{刚体运动}.
\end{definition}

\begin{theorem}\label{theorem:刚体运动等距变换的充要条件}
等距变换 $\sigma: E^3 \to E^3$的充要条件是存在常向量$\bm a\in E^3$和$3$阶正交矩阵 $A$ (即$A\in O(3)$),使得
\begin{align}
\sigma(\bm x) = \bm a + A\bm x, \qquad \bm x\in E^3. \label{eq:isometry}
\end{align}
进一步, $\sigma: E^3 \to E^3$是刚体运动当且仅当 $\det A = 1$ .存在常向量$\bm a\in E^3$,以及$3$阶正交矩阵 $A$且$\mathrm{det}A=1$(即 $A\in SO(3)$),使得
\begin{align*}
\sigma(\bm x) = \bm a + A\bm x, \qquad \bm x\in E^3.
\end{align*}
\end{theorem}
\begin{proof}

充分性都显然成立.下证必要性.
在 $E^3$ 中取定一个正交标架 $\{O;\bm i,\bm j,\bm k\}$, 并将点与它的坐标等同. 设 $\sigma$ 保持距离, 则 $\sigma$ 将单位正交标架映成单位正交标架, 且保持内积. 令 $\bm a = \sigma(\bm 0)$, 并记 $A$ 为 $\sigma$ 在标准基下的线性部分, 则 $A$ 为正交矩阵, 且 $\sigma(\bm x)=\bm a+A\bm x$. 反之, 形如 \eqref{eq:isometry} 的变换显然保持距离. 若 $\sigma$ 还保持右手系, 则 $\det A = 1$; 若 $\det A = -1$, 则 $\sigma$ 是一个刚体运动与一个反射的合成.
\end{proof}

\begin{theorem}
设 $\{O;\bm i,\bm j,\bm k\}$ 和 $\{p;\bm e_1,\bm e_2,\bm e_3\}$ 是 $E^3$ 中任意两个右手正交标架, 则存在唯一的刚体运动 $\sigma$, 使得
\begin{align*}
\sigma (O)=p\text{且}\sigma (\boldsymbol{i})=\boldsymbol{e}_1,\sigma (\boldsymbol{j})=\boldsymbol{e}_2,\sigma (\boldsymbol{k})=\boldsymbol{e}_3.
\end{align*}
\end{theorem}
\begin{remark}
刚体运动可以解释为坐标变换: 若点 $q$ 在标架 $\{O;\bm i,\bm j,\bm k\}$ 下的坐标为 $(x,y,z)$, 在标架 $\{p;\bm e_1,\bm e_2,\bm e_3\}$ 下的坐标为 $(\tilde x,\tilde y,\tilde z)$, 则由这个定理知,存在唯一的将标架 $\{O;\bm i,\bm j,\bm k\}$映到标架 $\{p;\bm e_1,\bm e_2,\bm e_3\}$的刚体运动$\sigma$.由\refthe{theorem:刚体运动等距变换的充要条件}知存在$\bm a\in E^3$和$A\in SO(3)$,使得
\begin{align*}
\sigma(\bm x) = \bm a + A\bm x, \qquad \bm x\in E^3.
\end{align*}
进而
\begin{align*}
(x,y,z) = \bm a + (\tilde x,\tilde y,\tilde z)A,
\end{align*}
\end{remark}
\begin{proof}

设 $\bm a = \overrightarrow{Op}$, $A$ 是由 $(\bm e_1,\bm e_2,\bm e_3)$ 在标架 $\{O;\bm i,\bm j,\bm k\}$ 下的坐标构成的正交矩阵, 其行列式为 $1$. 定义 $\sigma(\bm x)=\bm a + A\bm x$, 则 $\sigma$ 是刚体运动且满足要求. 唯一性由刚体运动在一点和一组基上的作用完全决定.
\end{proof}

\begin{proposition}
$E^3$ 中全体刚体运动的集合与 $E^3\times SO(3)$ 一一对应, 因而构成一个 $6$ 维流形.
\end{proposition}
\begin{proof}



\end{proof}

\begin{definition}[仿射标架]
\textbf{仿射标架}定义为$E^3$ 中一点 $p$ 与三个不共面的向量 $\bm e_1,\bm e_2,\bm e_3$ 组成的图形 $\{p;\bm e_1,\bm e_2,\bm e_3\}$, 不要求单位正交. 其\textbf{度量系数}定义为
\begin{align*}
g_{ij} = \bm e_i\cdot\bm e_j,\qquad 1\leqslant i,j\leqslant 3.
\end{align*}
\end{definition}

\begin{proposition}
全体仿射标架的集合与 $E^3\times GL(3)$ 一一对应, 构成一个 $12$ 维空间.
\end{proposition}
\begin{proof}



\end{proof}





\end{document}